\section{Diagnosis in detail: the plausibility ceiling, the confusion matrix, and the
reordering experiment}\label{sec:D}

\subsection{The monitor's eighteen named defaults, and what the plausibility ceiling costs}\label{sec:D.1}

Every one of those numbers is a named field with a documented default rather than a literal buried
in a branch, which is what let the register of §\ref{sec:11.6} name what each experiment must re-measure:
$k$ = 30 ppm, $h$ = 150 ppm, cycle 60 s, suspect fraction 0.6, confirm persistence 2 cycles,
suspect-clear 3 cycles, warmup 3 cycles, stale 3 intervals, evaluation window 10 min, generation
gain 0.3, anchor gain 0.3, minimum step 0.05, probe interval 30 min, probe step $+0.20$,
common-mode 5 of 6, step-amplitude match 0.3, overnight tolerance 100 ppm over 24 h, integration
step 1 s. One of these deserves an argument rather than a line in a list: the plausibility ceiling
runs at \textbf{1.25} times design occupancy, and 25 \% of headroom above design is a deliberately
desensitising choice. It is now priced rather than asserted.

\texttt{MonitorParameters.withOccupancyCeiling} lets T-24 sweep the headroom against a control arm
with nothing injected, and the sweep settles what the setting costs: at a ceiling of 1.00 the
closed-loop arm detects \textbf{16 of 20} at a median of 126 min from onset, but latches a room in
\textbf{2 of 20} fault-free runs; at the shipped 1.25 it detects \textbf{15 of 20} at a median of
180 min from onset and latches \textbf{0 of 20}. Those minutes are counted from fault onset so
that the two ceilings are comparable, and they are not R-1's delay, which Section \ref{sec:11.2}'s
closed-loop arm counts from the first ignored command and reports as a median of 74 min and a p95
of \textbf{257 min}. Tightening the ceiling therefore buys one detection and
some speed, and pays two false alarms for them, so the headroom is not a free lever between R-1
and R-3. And below 1.00 it stops bounding occupancy at all.

That control arm also retires a hypothesis this report would otherwise have left standing. A
closed-loop detection rate quoted with nothing to compare it against is a count, not a
measurement, and it was reasonable to suspect that the 15 detections were R-3's false-alarm
mechanism landing on a room that happened to have a fault in it. The fault-free column is 0 of 20
at the shipped headroom, so they are real detections. The arithmetic behind the sweep is the same
one that limits R-1 in Section \ref{sec:13.4}: at A124's drawn occupancy of four to five of its eight desks
the frozen damper is worth about \textbf{86 ppm} of steady state, against a ceiling sitting
\textbf{270--600 ppm} above the reading.

\subsection{What the classifier actually emits, measured (T-15)}\label{sec:D.2}

The four-signal table of §\ref{sec:7.3} is what the tree is \emph{for}; this is what it
does. \texttt{ConfusionMatrixExperiment} runs the pipeline once per
injected condition (the four injectable faults plus a no-fault arm) on a fixed seed, and
tabulates the condition injected against the classification emitted. The no-fault row is why T-15
is traced to R-3 and not only to R-1. 218 runs in all: 110 at the headline magnitudes over two arms,
plus a 54-run drift-rate sweep and a 54-run offset sweep. Each cell is the first confirmed case
anywhere in the building, because that is the work order an operator actually receives. Note that
the rows are \emph{mechanisms} and the columns are \emph{classifications}, and there are more of the
first than the second: \texttt{SENSOR\_DRIFT} and \texttt{SENSOR\_OFFSET} are two ways for a sensor
to be wrong and the monitor has one word for both.

\begin{footnotesize}
\begin{verbatim}
seed 42, closed loop   STUCK_DAMPER  BLOCKED_FILTER  SENSOR_DRIFT  UNCLASSIFIED  never confirmed
nothing injected        .             .               .             .             1
STUCK_DAMPER            1             .               .             .             .
BLOCKED_FILTER          .             .               .             1             .
SENSOR_DRIFT            1             .               .             .             .
SENSOR_OFFSET           1             .               .             .             .
\end{verbatim}
\end{footnotesize}

Over seeds 42 and 1--10, eleven runs per row: nothing injected is silent in 10 of 11,
\texttt{STUCK\_DAMPER} is named correctly in 8 of 11 and never confirms in 3, \texttt{BLOCKED\_FILTER}
comes out \texttt{UNCLASSIFIED} in 11 of 11, \texttt{SENSOR\_DRIFT} comes out \texttt{STUCK\_DAMPER}
in 11 of 11, and \texttt{SENSOR\_OFFSET} comes out \texttt{STUCK\_DAMPER} in 4, \texttt{UNCLASSIFIED}
in 6 and never confirms in 1. \textbf{One of the four mechanisms is classified correctly}, and the
report says so here rather than leaving the tree's branches to imply otherwise. Three things are off
the diagonal and each has a mechanism:

\begin{itemize}
\item \textbf{\texttt{SENSOR\_DRIFT} is never emitted: 0 of 76 runs from a ramp and 0 of 76 from a
step offset}, across rates 5--160 ppm/h and offsets 50--700 ppm, over both arms. The leaf asks for
two things in the same instant: an informative overnight floor \emph{and} a step response that still
matched. \textbf{The two injectable sensor mechanisms each supply one of them and never both.} A
ramp outlives the 24 h overnight window (it keeps growing, so it is still confirmable after a day),
and by then the offset it accumulated has inflated the quasi-steady generation estimate past the
step amplitude the leaf needs. A step offset keeps its amplitude and never reaches the window: with
the inner loop running it confirms within the hour, and with it off nothing re-anchors the observer,
the quasi-steady estimate absorbs a constant, and it never confirms at all. Counted over every
sensor-fault run that confirmed, the ramp reached an informative floor 12 times and a matched step 0
times, the step reached a matched step 15 times and an informative floor 0 times, and
\textbf{neither reached both once}.

This report previously called that \emph{a gap in the fault model, not a defect in the classifier},
on the grounds that the classifier's own unit test reaches the branch through a double applying a
constant offset instantaneously, which no ramp can. \textbf{That attribution was tested and it was
wrong.} The generator was given exactly that shape (\texttt{FaultKind.SENSOR\_OFFSET}, injectable
through the unchanged four-field \texttt{FAULT\_SPEC} grammar), and the branch is still unreachable.
The binding constraint is that \textbf{the classification is computed at the confirmation while one
of its inputs needs a day of the fault}, which is a property of when the tree runs. The arithmetic
that closes it: the overnight floor is a minimum over 1440 noisy samples, so it sits about 3.3\,$\sigma$
(very nearly 100 ppm) below the level the room is holding, and the closest any offset came was
a floor of 489 ppm against the 520 ppm the leaf needs, on a case confirmed 1621 min in. An offset
large enough to lift the floor past 520 confirms within three minutes.
\item \textbf{\texttt{BLOCKED\_FILTER} confirms as \texttt{UNCLASSIFIED}}, 11 of 11, exactly as
Section \ref{sec:4}'s cross-room analysis predicts: the proportional split only redistributes a fixed total
flow, so no command the monitor can watch pushes all six residuals the same way. In every run the
first room confirmed is a room with \emph{no} damper fault in it, and all six latch by the end of
the window. The fan-power excess is still attached as evidence, so the work order points at the
plant even though the classification does not.
\item \textbf{The stuck-damper leaf fires with no step on record, and trying to stop it emptied
the leaf.} This is the most instructive thing T-15 found, and it took a failed fix to see it.
\texttt{stepShapeUnchanged()} used to return \texttt{false} when the step-amplitude ratio was
\texttt{NaN}, and \texttt{NaN} means \emph{no step was ever judged}. So the test ``the commanded
step produced no matching response'' was satisfied by a room that had never been commanded to move.
9 of 57 drift runs confirm by that route with zero judged steps; the other 48 failed the amplitude
test genuinely, on 1--419 steps. The same codebase gets this right one field away, where
\texttt{FaultSignals} documents the overnight floor's \texttt{NaN} as ``no evidence, not healthy''.

The signal is now a three-valued \texttt{StepResponse} (\texttt{MATCHED},
\texttt{MISSING}, \texttt{NOT\_JUDGED}), so the conflation cannot be written by accident any
more. \textbf{But the leaf still accepts \texttt{NOT\_JUDGED}, deliberately, and the reason is a
property of the architecture rather than a shortcut.} A probe is requested \emph{through the
evidence on} \texttt{FaultStateChanged}, so a case must already be open before any step is
commanded, which means \textbf{at the moment of classification no step has been judged, ever}.
Requiring positive evidence of an unanswered step was tried and measured: it took the one cell on
T-15's diagonal to \texttt{UNCLASSIFIED} and left the tree classifying \emph{nothing at all}:
the one named mechanism becoming none. The step-amplitude signal arrives \emph{after} the decision it was
designed to inform, and no amount of care in the classifier can fix an ordering problem in the
pipeline.

So the leaf fires on localisation and flat fan power, and the work order now \textbf{says so}: a
diagnosis reached without a judged step carries an explicit entry that it rests on two signals
rather than three, that the room has not been asked to move and therefore has not failed to answer,
and that a probe has been requested so the third signal can confirm or overturn it. The weakness is
unchanged in behaviour and no longer invisible, which is what the evidence mechanism is for, and
is strictly better than a boolean that quietly asserted a test nobody ran.

\textbf{The obvious repair was then built and measured, and it does not work, because the third
signal is blind after a confirmation rather than merely late.} If the problem is ordering, reorder:
latch the case on the same cycle as before but \emph{withhold the diagnosis}, request the probe on
that cycle instead of one probe interval later, and name the fault when the step window closes. That
is implemented (\texttt{ClassifyWhen.AFTER\_PROBE}) and measured as T-35 over 176 paired runs on
identical seeds. It does what it was designed to do: the case latches on the same cycle (detection
moves in \textbf{0 of 88 pairs}), and the fault is named on a genuinely judged step in 28 of the 39
runs that defer, against 0 under the deployed order. \textbf{And the settled classification is
identical in all 88 pairs}, because the judged step comes back \texttt{MISSING} for a stuck damper,
for a blocked filter, for nine drift runs whose damper is in perfect health, \emph{and} for the run
with nothing injected at all. A signal that reads the same for every condition decides nothing, so
the reorder buys a median 11 min of work order with no diagnosis on it and no change of answer. It
is \textbf{not adopted}, and \texttt{ON\_CONFIRMATION} remains the deployed ordering.

The cause is located and it predates the ordering. \texttt{ZoneObserver.recordStepAmplitude}
accumulates the measured leg of the amplitude against \texttt{predictedAtStep} while the predicted
leg is also against \texttt{predictedAtStep}, so the measured leg carries the \textbf{standing
residual} as a constant offset. The anchor-gain pull that would make that baseline legitimate is
switched off while a case is open, so at the confirmation the prediction is stale by exactly the
residual that confirmed it, and the ratio becomes a statement about the residual rather than about
the response. On a 600 ppm offset with a healthy damper that answers its probe by dropping 300 ppm,
the ratio runs from $-7.0$ to $-0.8$: negative. \texttt{ZoneObserver.measuredAtStep} is assigned for
this purpose and never read.

Correcting that baseline was measured too, as an unlanded counterfactual, and it makes the reorder
\emph{worse}: 18/44 to 15/44 on T-15's diagonal, because the corrected signal reads \texttt{MATCHED}
on 3 of 8 genuinely stuck dampers and \texttt{MISSING} on 5 of 11 healthy ones. By the time a probe
can be commanded the supervisor has degraded the zone, raised the fan and trimmed the five healthy
dampers, so \textbf{the room's answer to the probe is mostly the compensation's answer rather than
the damper's}. The 18/44 baseline reproduces on the current code; the 15/44 was measured on a one-line patch that was then reverted, so it is the one figure in this report that cannot be re-derived from the repository as it stands. \textbf{The conclusion is the useful part}: a three-signal stuck-damper diagnosis
needs a step judged \emph{before} a case opens (the inner loop's own trims, read against a
baseline that survives them), not a probe requested after the response to the fault has already
begun. That is design work and §\ref{sec:13} carries it.
\end{itemize}

\subsection{Why the sensor-drift branch is the classifier's and not the generator's}\label{sec:D.3}

A second, sharper limit on what can be claimed, and T-15 has since made it sharper still: of the
faults the project set out to name, \textbf{one is classified and the others can only be
pointed at}. This paragraph said \emph{two} until the confusion matrix was measured, and the
correction is the single largest change the measurement made to this report.
\texttt{STUCK\_DAMPER} is the reliable one, and T-24, T-25, T-26 and T-34b are all measurements of
it. \texttt{BLOCKED\_FILTER}'s branch is \textbf{unreachable with this physics} for structural
reasons that Section \ref{sec:13.1} sets out: a blocked filter confirms as \texttt{UNCLASSIFIED} with its
plant power excess reported, which is the right behaviour for a case the system cannot prove.
\texttt{SENSOR\_DRIFT} was described here as classified \emph{conditionally}, on the grounds that
its signature needs 24 h of overnight floor and a case opened earlier comes out
\texttt{UNCLASSIFIED}. Measured, the condition is not merely restrictive but \textbf{unreachable}: the branch was emitted
in \textbf{0 of 76 runs} across drift rates from 5 to 160 ppm/h, because the injected drift is a
\emph{rate} and by the time the floor has become informative the accumulated offset has already
pushed the step amplitude below its match fraction. The room is then named a stuck damper: a
confident, specific, wrong work order rather than an honest \texttt{UNCLASSIFIED}.

\textbf{That diagnosis was then tested, and it was wrong in an instructive direction.} The reasoning
above blames the \emph{fault model}: the classifier's own unit test reaches the branch through a
double that applies a constant offset instantaneously, which no ramp can produce, so the obvious
remedy was to give the generator that shape. It now has it: \texttt{FaultKind.SENSOR\_OFFSET}, a
sensor whose zero jumps once and holds, which is a real failure mode in its own right (a
recalibration against the wrong reference, a replacement element installed without one) and is
injectable through the unchanged \texttt{KIND:TARGET:MAGNITUDE:ONSET} grammar. \textbf{The branch is
still unreachable}: 0 of 76 runs from a step offset, swept from 50 to 700 ppm, beside the 0 of 76
from a ramp. The leaf wants an informative overnight floor \emph{and} a matched step response in the
same instant, and the two mechanisms fail it in complementary places: of the runs that confirmed,
the ramp reached an informative floor 12 times and a matched step 0 times, the step reached a matched
step 15 times and an informative floor 0 times, and \textbf{neither reached both once}. The binding
constraint is that the classification is computed at the confirmation while one of its inputs needs a
day of the fault. \textbf{So this is a property of the classifier, not of the generator}, and the
report had the attribution the wrong way round until the experiment was run.

\textbf{The confusion matrix this paragraph said must be read with those caveats in front
of it now exists}, in Section \ref{sec:7.3}, and it is read with exactly them.
